\documentclass[10pt,a4paper]{amsart}
\usepackage[margin=19mm]{geometry}
\usepackage{amsmath,amssymb,mathtools}
\usepackage[scr=rsfs,cal=euler]{mathalfa}
\usepackage{xfrac}
\usepackage[unicode,hidelinks,bookmarks=false]{hyperref}
\newcommand{\ie}{i.e.\ }
\newcommand{\cf}{cf.\ }
\newcommand{\resp}{respectively\ }
\newcommand{\Subsection}{\subsection}
\providecommand{\nobreakdash}{\nobreak\hbox{-}}
\setlength{\emergencystretch}{2em}
\multlinegap0pt
\newcommand{\mres}{\mathbin{\vrule height 1.6ex depth 0pt width
0.13ex\vrule height 0.13ex depth 0pt width 1.3ex}}
\newtheorem{thm}{Theorem}[section]
\newtheorem{cor}[thm]{Corollary}
\newtheorem{claim}[thm]{Claim}
\newtheorem{lemma}[thm]{Lemma}
\newtheorem{theorem}{Theorem}
\newtheorem{conjecture}{Conjecture}
\setlength{\topmargin}{0in}
\newtheorem{example}{Example}
\newtheorem{rmk}[thm]{Remark}
\newtheorem{problem}[thm]{Problem}
\newtheorem{defn}[thm]{Definition}
\newtheorem{prop}[thm]{Proposition}
\newtheorem{conj}[thm]{Conjecture}
\newtheorem{question}[thm]{Question}
\newtheorem{ack}[thm]{Acknowledgements}
\numberwithin{equation}{section}
\begin{document}
\title[Minimal equatorial fillings]{Minimal equatorial fillings}
\author{Jacob Bernstein; Daniel Ketover}
\date{}
\maketitle
\noindent\emph{Source and licence.} Minimal equatorial fillings. Version arXiv:2609.26701v1 (CC BY 4.0). This material is distributed under the Creative Commons Attribution 4.0 International licence, http://creativecommons.org/licenses/by/4.0/, as recorded in the arXiv OAI licence metadata for this exact version and shown on the abstract page https://arxiv.org/abs/2609.26701v1. Full rights notice: licenses/arxiv-2609.26701-CC-BY-4.0.txt.\par\smallskip
\noindent\textit{Editorial scope.} Counted unit: the original \texttt{thm} environment \texttt{admissible} at source lines 810--819 together with its complete proof at lines 822--870; the delivered document keeps source lines 783--871 so that every referenced label and every citation resolves inside it. Native editable source is the paper's own concatenated TeX; the AMS wrapper is editorial formatting only. No publisher class, upstream script or unrelated artwork is redistributed.
\section{Admissible genera}\label{admissiblesection}
In this section we collect results of Whitney and Massey that give the possible combinations of genera and Euler numbers can be realized by smooth embeddings with boundary a great circle. 

If $\Sigma$ is a non-orientable embedded surface of genus $g$ in $\mathbb{S}^3$ with boundary an unknot $K$, then capping it off with a disk in $B^4$ we obtain a closed non-orientable surface $\tilde{\Sigma}$ embedded in $\mathbb{R}^4$ with  $\chi(\tilde{\Sigma})=2-g$.  

Whitney \cite{Whitney} proved that the normal Euler number $e(\tilde{\Sigma})=e(\Sigma)$ (see Section 2 in \cite{Conway} for this equality or \cite[1.1.1]{Yasuhara}) and Euler characteristic $\chi(\tilde{\Sigma})$ obey the following relation 
\begin{equation}
e(\tilde{\Sigma})=2\chi(\tilde{\Sigma}) \;\;\;\mbox{   mod } 4, 
\end{equation}
which gives
\begin{equation}
e(\Sigma)=-2g \;\;\;\mbox{   mod } 4.
\end{equation}


We deduce the following:

\begin{lemma}\label{parity}
If $\Sigma\subset \mathbb{S}^3$ is a non-orientable embedded surface with boundary an unknot with $e(\Sigma)=2k$, then the genus of $\Sigma$ has the same parity as the integer $k$.
\end{lemma}

Massey \cite{Ma} -- see also \cite[Corollary 1.1.1]{Yasuhara} -- later proved that 
\begin{equation}\label{options}
e(\Sigma)\in \{-2g,-2g+4,..., 2g-4, 2g\}.
\end{equation}

We obtain: 

\begin{thm}[Admissible genera \cite{Whitney,Ma}]\label{admissible}
There exists a smooth embedded non-orientable surface $\Sigma$ with boundary an unknot in $\mathbb{S}^3$ if and only if
\begin{equation}\label{adgen}
\mbox{genus}(\Sigma)=\frac{e(\Sigma)}{2}+2n\mbox{ where } n\in\mathbb{N}\cup\{0\}.
\end{equation}
In particular for such a surface $\Sigma$,
\begin{equation}
\mbox{genus}(\Sigma)\geq\frac{e(\Sigma)}{2}.
\end{equation}
\end{thm}

\begin{proof}
Lemma \ref{parity} and \eqref{options} establish the necessity of \eqref{adgen} for such a surface $\Sigma$ to exist.  Given $k\geq 1$, we will show on the other hand that there exists a smooth embedded surface of genus $k$ and Euler number $2k$ and boundary an unknot.  Since one can then add $n$ handles to this surface (which increases the genus by $2n$ but does not change the Euler number), this shows the sufficiency of \eqref{adgen}.

Let us construct the desired surface in our setting.  Fix $k\geq 1$. Consider the immersed band $\overline{\tau}_{1,2k}$ restricted to a neighborhood $T_{\pi/4}(C^*)$ of the curve $C^*$ where it fails to be an embedding.  The solid torus $T_{\pi/4}(C^*)$ is foliated by the meridian disks $\{H_\theta\cap  T_{\pi/4}(C^*)\}_{t\in [0,2\pi]}$.  The cross sections \begin{equation}F_\theta = H_\theta\cap T_{\pi/4}(C^*)\cap \overline{\tau}_{1,2k}\end{equation} consist of $2k$ equally spaced geodesic segments meeting at the center point $H_\theta\cap C^*$ of $H_\theta$.  

For $\delta$ small enough, let $B_\delta\subset H_0$ be the geodesically convex ball of radius $\delta$ in $H_0$ based at its center.  Then 
\begin{equation}
(\bigcup_{j=1}^{2k} \gamma_{0,\frac{2\pi j}{k}})\cap \partial B_\delta = p_1(\delta)\cup...\cup p_{2k}(\delta), 
\end{equation}
where $p_1(\delta),...,p_{2k}(\delta)$ are equally spaced points on $\partial B_\delta$ varying continuously with $\delta$.
Let $g_i(\delta)$ denote the unique minimizing geodesic segment in $H_0$ between $p_i(\delta)$ and $p_{i+1}(\delta)$ (where $g_{2k}(\delta)$ is the unique minimizing geodesic between $p_{2k}(\delta)$ and $p_1(\delta)$).  Note that the interiors of the segments $\{g_i(\delta)\}_{i=1}^{2k}$ are pairwise disjoint since they are minimizing, the segments vary continuously in $\delta$, and are contained in $B_\delta$.

Let us denote the piecewise smooth geodesic contour:
\begin{equation}
G(\delta) = \bigcup_{i=1}^{2k} g_i(\delta).
\end{equation}


For $\epsilon<\frac{\pi}{4}$, let  $\{\Psi_t\}_{t\in [0,\epsilon]}\subset F_0$ be the family of curves given by 
\begin{equation}
\Psi_t = G(t)\cup \bigcup_{j=1}^{2k} \gamma_{0,\frac{2\pi j}{k}}\cap (F_0\setminus B_t) 
\end{equation}

Let us extend $\{\Psi_t\}_{t\in [0,\epsilon]}$ to a family $\{\Gamma_t\}_{t\in [0,2\pi]}$ also contained in the disk $F_0$ by setting
\begin{equation}
\Gamma_t = 
\begin{cases} 
  \Psi_t & \text{if } t\in [0,\epsilon]\\ 
  
  \Psi_\epsilon& \text{if } t\in [\epsilon,2\pi-\epsilon]\\
  \Psi_{2\pi-t} & \text{if } t\in [2\pi-\epsilon,2\pi]. \\
\end{cases}
\end{equation}
Note that for all $t\in [0,2\pi]$
\begin{equation}\label{welldefined}
\Gamma_t\cap \partial T_{\pi/4}(C^*) = F_0\cap \partial T_{\pi/4}(C^*) =\overline{\tau}_{1,2k}\cap \partial T_{\pi/4}(C^*).
\end{equation} Then define an embedded (piecewise smooth) surface by setting
\begin{equation}
S_k=\bigcup_{\theta\in [0,2\pi]}\sigma_\theta \circ \rho_{\frac{\theta}{2k}}(\Gamma_\theta).
\end{equation}
Roughly speaking, the cross sections $S_k\cap F_\theta$ first desingularize the $2k$ segments comprising $F_0$ in one of the two possible ways, then rotate the result around $C^*$ so that it comes back with a $\frac{\pi}{k}$ twist desingularizing $F_0$ in the other possible way. 

Finally we define the desired surface with boundary $C$ (in light of \eqref{welldefined}) by:
\begin{equation}
G_k = (\overline{\tau}_{1,k}\cap (\mathbb{S}^3\setminus T_{\pi/4}(C^*)))\cup S_k.
\end{equation}
Since $S_k$ agrees with $\overline{\tau}_{1,2k}$ in a neighborhood of $C$ it follows that \begin{equation}e(S_k)=e(\overline{\tau}_{1,2k})=2k.\end{equation}  
One can readily compute that the genus of $S_k$ is $k$.
\end{proof}

\begin{thebibliography}{99}
\bibitem[Conway]{Conway} A.~Conway, P.~Orson, and M.~Powell, \emph{Unknotting nonorientable surfaces}, to appear JEMS (2023).
\bibitem[Ma]{Ma} W.~S. Massey, \emph{Proof of a conjecture of {W}hitney}, Pacific J. Math. \textbf{31} (1969), 143--156.
\bibitem[Whitney]{Whitney} H.~Whitney, \emph{On the topology of differential manifolds}, Michigan University Press, 1940.
\bibitem[Yasuhara]{Yasuhara} A.~Yasuhara, \emph{Connecting lemmas and representing homology classes of simply connected 4-manifolds}, Tokyo Journal of Mathematics \textbf{19} (1996), no.~1.
\end{thebibliography}
\end{document}
